% 本文件由 LaTeX 排版 Agent 根据有效 translation.json 自动生成。
% author=Julius Africanus; work_id=0614; title=现存著作（尤利乌斯·阿非利加努斯）

\chapter{现存著作（\Person{尤利乌斯·阿非利加努斯}）}

% structure: p0001 source=h1 target=omitted reason=page-title-already-rendered-by-parent

% structure: p0002 source=h2 target=section reason=relative-heading-rank; base=h2

\section{致\Person{亚里斯提德}书}

% structure: p0003 source=h3 target=subsection reason=relative-heading-rank; base=h2

\subsection{一}

……有些人确实错误地声称，这种对司祭和君王之名的不同列举和混合——正如他们所想的那样——是适当进行的，为的是表明基督有权利同时是司祭和君王；好像有人不相信这一点，或者有别样的希望，而不相信基督是他圣父的大司祭，将我们的祈祷呈献给祂，又是超世间的君王，以圣神治理他所拯救的人，是万物治理的合作者。这一点不是通过支派的目录，也不是通过记载的世系的混合，而是通过宗祖和先知向我们宣布的。因此，我们不要堕落到如此宗教性的琐碎中，即通过名字的互换来确立基督的君王职和司祭职。因为肋未的司祭支派也与犹大的君王支派联姻，通过\Person{亚郎}娶了\Person{纳赫雄}的姐妹\Person{依撒伯尔}，以及\Person{厄肋阿匝尔}又娶了\Person{法提耳}的女儿并生儿育女。因此，福音作者们就会说了假话，肯定了不真实的东西，而非真实的赞美。为此，一位追溯了\Person{若瑟}的父亲\Person{雅各伯}的谱系，从\Person{达味}经过\Person{撒罗满}；另一位则追溯了\Person{若瑟}的父亲\Person{赫里}的谱系，虽然方式不同，从\Person{达味}的儿子\Person{纳堂}开始。他们本不应不知道，所列举的两个祖先世系都是\Person{达味}的后裔，即犹大的君王支派。因为如果\Person{纳堂}是先知，那么\Person{撒罗满}也是，他们两人的父亲也是；许多支派中都有先知，但除了肋未支派外，没有一个支派有司祭。因此，他们的捏造毫无意义。这种主张也不会在基督的教会中胜过精确的真理，以至于为了基督的赞美与荣耀而编造谎言。因为谁不知道那位宗徒的至圣之言呢？当他宣讲并宣告我们救主的复活，并确信地坚持真理时，他极其敬畏地说：\Quote{若有人说基督没有复活，而我们却断言并相信了这事，且盼望并传扬这事，我们就是天主的假见证，因为声称天主复活了基督，而实际上他没有复活。} 如果这位光荣天主圣父的人如此害怕自己在叙述一件奇妙事实时被视为假见证，那么那个试图用虚假陈述来确立真理、制造错误观点的人，岂不更应害怕吗？因为如果谱系不同，且没有追溯到\Person{若瑟}的真正后代，而所有陈述只是为了确立那将要诞生者的地位——即确认真理，即那将要来的将是君王和司祭——同时没有提供任何证明，反而使言语的尊严沦为一首虚弱的赞美诗——那么显然，天主不会因此得到任何赞美，因为它是一个谎言，反而审判会归于那断言者，因为他把虚假当作真实来夸耀。因此，为了揭露那如此说话之人的无知，并防止任何人因这愚行而跌倒，我将陈述这些事的历史真相。\par

% structure: p0005 source=h3 target=subsection reason=relative-heading-rank; base=h2

\subsection{二}

因为在以色列中，他们世代的名字要么按本性、要么按法律被列举——按本性，确实是通过合法后裔的继承；按法律，则是当另一个人为无子女而死的兄弟名字兴起子孙时；因为那时还没有给予他们复活的明确希望，他们以一种可朽的复活形式拥有将来应许的表征，目的是为了延续死者的名字。因此，在这个谱系中记载的人中，有些是按合法继承作为儿子继承父亲，而另一些是在一个家族中出生但被转移到另一个家族名下，所以二者都被提及——那些实际是祖先的人，以及那些只是名义上如此的人。因此，两位福音作者都没有错误，因为一位按本性计算，另一位按法律计算。因为几个世代，即从\Person{撒罗满}下来的和从\Person{纳堂}下来的，如此交织在一起——通过为无子女者兴起孩子，再婚，以及兴起后裔——以至于同一个人往往完全公正地被认为时而属于这一个，时而属于另一个，即属于他们名义上的父亲或实际的父亲。因此，这两个记述都是真实的，并且一直追溯到\Person{若瑟}，虽然相当复杂，但却非常准确。\par

% structure: p0007 source=h3 target=subsection reason=relative-heading-rank; base=h2

\subsection{三}

为使以上所述显明，我将解释世代之间的替换。若我们从达味经撒罗满计算世代，\Person{玛堂}正是倒数的第三位，他生了\Person{雅各伯}，即\Person{若瑟}的父亲。但若我们随\Person{路加}从\Person{达味}之子\Person{纳堂}计算，则倒数的第三位是\Person{默尔希}，他的儿子是\Person{赫里}，即\Person{若瑟}的父亲。因为\Person{若瑟}是\Person{赫里}的儿子，\Person{赫里}是\Person{默尔希}的儿子。既然\Person{若瑟}是我们所讨论的对象，我们必须表明何以两方都被视为他的父亲——既说\Person{雅各伯}出于\Person{撒罗满}，又说\Person{赫里}出于\Person{纳堂}：首先，\Person{雅各伯}和\Person{赫里}这二人如何是兄弟；其次，这两人的父亲\Person{玛堂}和\Person{默尔希}虽属不同家族，却如何被表明是\Person{若瑟}的祖父。那么，\Person{玛堂}和\Person{默尔希}相继娶了同一女子为妻，所生的孩子是同母异父的兄弟，因为法律并不禁止寡妇——无论是因离婚或因丈夫死亡——再嫁他人。于是，由\Person{厄斯塔}——根据传承，这是她的名字——\Person{玛堂}先是\Person{撒罗满}的后裔，生了\Person{雅各伯}；\Person{玛堂}死后，\Person{默尔希}——他的世系追溯至\Person{纳堂}——同属一支派但不同家族，正如已说娶了她，生了一个儿子\Person{赫里}。这样，我们将发现\Person{雅各伯}和\Person{赫里}是同母异父的兄弟，尽管家族不同。其中，\Person{雅各伯}娶了他兄弟\Person{赫里}的妻子——\Person{赫里}无子而死——由她生了第三个孩子\Person{若瑟}——按本性和按记录都是他的儿子。因此也记载说：\BibleQuote{雅各伯生若瑟}。但按法律，他是\Person{赫里}的儿子，因为他的兄弟\Person{雅各伯}为他立嗣。因此，通过他所推导的族谱也不会无效，福音史家\Person{玛窦}在他的列表中是如此给出的：\BibleQuote{雅各伯生若瑟}。但\Person{路加}却说：\BibleQuote{若瑟——按人看是\Person{赫里}的儿子、\Person{默尔希}的孙子}（因为他自己也加上了这一句）。因为不可能更明确地陈述按法律的世系；因此在这种生成方式中，他完全省略了\Quote{生}这个词直至最后，将族谱以结论方式追溯至\Person{亚当}和\Person{天主}。\par

% structure: p0009 source=h3 target=subsection reason=relative-heading-rank; base=h2

\subsection{四}

这确实并非无法证明，也并非轻率的猜测。因为救主按血肉的亲属，无论为夸耀自己的出身或只是陈述事实，总之是说实话，也传下了以下记述：一些厄东盗匪袭击\Place{巴勒斯坦}的\Place{阿斯卡隆}城，除了从建在城墙附近的一座\Person{阿波罗}神庙掠走其他财物外，还掳走了\Person{安提帕特}——他是某位\Person{黑落德}的儿子，也是该庙的仆役。由于司祭无法为儿子支付赎金，\Person{安提帕特}便在厄东人的习俗中长大，后来与\Place{犹大}的大司祭\Person{希尔卡诺斯}结为好友。他受\Person{希尔卡诺斯}之托出使\Person{庞培}，并为其恢复了被其弟\Person{阿黎斯托步罗}所荒废的王国，因而得以获得\Place{巴勒斯坦}总督的头衔。当\Person{安提帕特}因其大幸遭嫉妒而被暗杀后，其子\Person{黑落德}继位，后来在\Person{安多尼}和\Person{奥古斯都}治下，由元老院下令被任命为\Place{犹大}王。他的儿子是\Person{黑落德}和其他分封侯。这些记述在希腊史籍中也有记载。\par

% structure: p0011 source=h3 target=subsection reason=relative-heading-rank; base=h2

\subsection{五}

但直到那时，希伯来人的族谱都登记在公共档案中，包括那些追溯到皈依者的——例如，\Person{阿希约尔}（阿孟人）、\Person{卢德}（摩阿布人），以及那些与以色列人一同离开埃及并与之通婚的人——\Person{黑落德}知道以色列人的谱系对他毫无帮助，且因意识到自己出身卑贱而受刺激，便焚烧了他们的家族记录。他这样做，是以为如果无人能通过公共档案将他的谱系追溯到先祖或皈依者，以及那称为\OriginalTerm{革奥勒}{georae}的混血种族，他就可以显得出身高贵。然而，少数勤奋之人拥有自己的私人记录，或是通过记忆名字，或是通过其他方式从档案中获取，他们自豪地保持着对高贵出身的记忆；其中恰好包括了前面提到的称为\OriginalTerm{得司波绪尼}{desposyni}的人，因他们与救主的家族有联系。这些人从\Place{纳匝拉}和\Place{科哈巴}——\Place{犹大}乡村——来到国家的其他地区，从\WorkTitle{编年纪}中尽可能准确地陈述上述族谱。那么，无论情况是否如此，根据我本人及任何明智者的意见，无人能发现更明显的解释。就以此事为足，尽管它没有见证支持，因为我们没有更令人满意或更真实的依据可提供。然而，福音在任何情况下都陈述真理。\par

% structure: p0013 source=h3 target=subsection reason=relative-heading-rank; base=h2

\subsection{六}

\Person{玛堂}是\Person{撒罗满}的后裔，生了\Person{雅各伯}。\Person{玛堂}死后，\Person{默尔希}是\Person{纳堂}的后裔，由同一妻子生了\Person{赫里}。因此\Person{赫里}和\Person{雅各伯}是同母异父的兄弟。\Person{赫里}无子而死，\Person{雅各伯}为他立嗣，生了\Person{若瑟}——按本性是自己的儿子，但按法律是\Person{赫里}的儿子。如此，\Person{若瑟}是两者的儿子。\par

% structure: p0015 source=h2 target=section reason=relative-heading-rank; base=h2

\section{基督诞生时波斯发生之事叙述}

关于这部作品的最佳导言，就是\Person{Migne}所刊的以下序言：\Quote{至今有许多博学之士认为，\Person{Africanus}关于基督诞生时波斯所发生之事的记述，是\Person{Sextus Julius Africanus}——一位基督后第三世纪的基督教学者——所著著名作品的片段；此作按时间顺序记述了直至\Person{Macrinus}在位期间的世界历史，分为五卷献给\Person{Mammaea}之子\Person{Alexander}，以期获准恢复其故乡\Place{Emmaus}。怀着与我所见激励\Person{Lambecius}及其提要作者\Person{Nesselius}同样的期望，我也以最大的热忱着手查阅我们选帝侯图书馆的手抄本。……但正如俗语所说，我找到的是煤炭而非宝藏。这段记述，远不能归诸一位为古代普遍口碑所誉的作家，其中毫无配得上编年史家\Person{Africanus}才华之处。因此，既然古代一致的见证表明他是一位学问渊博、判断力极强的人，而\WorkTitle{Cesti}的作者——也假托\Person{Africanus}之名——早已被批评家打上老妇轻信或惊人迷信倾向的烙印，我很容易赞同那些认为他与编年史家并非同一人，并将这篇拙劣之作也归于他名下的人的见解。但是，亲爱的读者，阅读这些篇章时，如果你的义愤不为作者的鲁莽所触动，至少你会与我一同嘲笑他惊人的愚行，同时也会明白，那些最杰出学问家的见证，不可估价过高以至于在条件许可时取代亲自的查验。}\par

% structure: p0017 source=h3 target=subsection reason=relative-heading-rank; base=h2

\subsection{波斯事件：论我主、天主及救主\Person{耶稣基督}的降生成人}

基督最初是从\Place{Persia}为人所知的。因为那个国家博学的律法家们，以最大的谨慎探究一切，没有什么能逃过他们的注意。因此，我将记录刻在金版上、保存在王家神庙中的事实；因为基督的名字正是从那里的神庙和与之相关的祭司听来的。那里有一座献给\Person{Juno}的神庙，甚至超过王宫，是那位受教于一切敬虔的王子\Person{Cyrus}所建，他在其中献上金银神像献给众神，并用宝石装饰它们——以免我浪费言辞详述那装饰。那时（正如金版记录所证），国王进入神庙，想要获得某些梦的解读，被祭司\Person{Prupupius}这样对他说道：\Quote{\InnerQuote{我恭喜你，主人：\Person{Juno}已经怀孕了。}国王微笑着对他说：\InnerQuote{那已死的人怀孕了吗？}他回答说：\InnerQuote{是的，那已死的又复活了，并且困扰生命。}国王说：\InnerQuote{这是什么？请解释给我听。}他回答说：\InnerQuote{说实话，主人，这些事的时候近了。}因为整夜，众男神和女神的雕像都在加热地面，彼此说：\InnerQuote{来，让我们恭喜\Person{Juno}。}他们对我说：\InnerQuote{先知，上前来；恭喜\Person{Juno}，因为她已被拥抱。}我说：\InnerQuote{那已不存在的人怎能被拥抱呢？}他们回答说：\InnerQuote{她又复活了，不再被称为\Person{Juno}，而是\Person{Urania}。因为大能者\Person{Sol}已经拥抱了她。}然后女神们对男神们说，把事情说得更明白：\InnerQuote{\Emph{Pege}就是被拥抱的那位；因为难道\Person{Juno}没有嫁给一位工匠吗？}男神们说：\InnerQuote{我们承认，她被称为\Emph{Pege}是正确的。她的名字，此外，是\Person{Myria}；因为她在子宫中，如同在深渊里，承载着一艘万他连得重量的船只。关于\Emph{Pege}这个称号，这样理解：这道水流发出永恒的灵之流——一道包含一条独一之鱼的流，这条鱼被神性的钩子钩住，并用它的肉供养整个世界，如同在海洋中一般。你说得好，她有一位工匠（作为配偶）；但那结合所生的工匠并不与她平等。因为这位降生的工匠，是首席工匠之子，以其卓越技艺架构了第三层天的穹顶，并以他的话语建立了这下界世界及其三重居住领域。}}\par

于是，雕像们就\Person{Juno}和\Emph{Pege}彼此争论，[最后]同声说：\Quote{当白昼结束，我们众神和女神都将清楚知道这事。所以，主人，请留到这一天结束。因为这事必将实现。因为所出现的事并非寻常。}\par

当国王留在那里观看雕像时，竖琴手们自动开始拨动竖琴，少女们歌唱；里面所有生物，无论四足动物或飞禽，用金银制成的，都发出了各自的声音。当国王战栗，充满极大恐惧时，他正要退去。因为他不能忍受这自发的大混乱。因此祭司对他说：\Quote{请留下，噢国王，因为完全的启示即将来临，是万神之神选择向我们宣告的。}\par

当这些话说完，屋顶打开了，一颗明亮的星降下，停在\Person{Pege}的圆柱上方，有一个声音这样说：\Quote{至高\Person{Pege}，大能之子派我来向你宣告，同时为您在分娩中服务，为要与您缔结无瑕的婚姻，噢，众生等级之首的母亲，三重神祇的新娘。由非凡生殖所生的婴孩被称作\Quote{开始}和\Quote{终结}——救恩的开始，与毁灭的终结。}\par

当这句话说出时，所有雕像都俯伏在地，只有\Person{Pege}的雕像站立着，其上还发现放置了一顶王冠，上侧有一颗镶嵌在红宝石和绿宝石中的星，下侧则靠着那颗星。\par

国王立即下令将他治下所有解梦者和智者带进来。当所有传令官迅速传令后，这些人都在神殿中聚集。他们看见\OriginalTerm{佩格}{Pege}上方的星、带有星和宝石的王冠，以及倒在地上的雕像，便说：\Quote{大王，一位神圣而君王的苗裔（后裔）已经兴起，带着天地之王的形象。因为佩格—米里亚是白冷人佩格的女儿。王冠是君王的标记，星是天上对地上将要发生之事的宣告。从犹大兴起了一个王国，将要颠覆犹太人所有的纪念物。众神倒伏在地，预示了他们尊荣的终结。因为那要来者，具有更古老的尊荣，将取代所有新近的。所以现在，大王，请派人去耶路撒冷。因为你将找到全能天主的基督，以身体形态降生在一个妇人的怀抱中。}星一直停留在被称为天上的佩格雕像上方，直到贤士们出来，然后它就跟随着他们。\par

然后，在深夜，狄俄倪索斯出现在神殿中，没有萨提尔陪伴，对众雕像说：\Quote{佩格不是我们中的一员，而是远在我们之上，因为她以神圣的方式孕育了一个人。祭司普鲁普皮乌斯！你还在这里耽搁什么？古时记载的一件行动已临到我们身上，我们将被一个充满力量和活力的人证伪。凡我们欺骗的，我们已欺骗；凡我们统治的，我们已统治。我们不再给神谕。我们的尊荣已离我们而去。我们变得没有荣耀，没有奖赏。只有一位，唯独一位，从所有人手中再次领受祂应有的尊荣。至于其余，不要不安。波斯人不再向天地索取贡物。因为建立这些的那一位就在眼前，要将实际的贡物带给派遣祂的那一位，更新古老的肖像，将肖像与肖像相合，使不相似变为相似。天与地同乐，地也因接受来自天的喜乐而欢跃。天上未曾发生的事，在低下之地发生了。蒙福秩序所未见的那一位，被悲惨的秩序所见。火焰威胁那些，甘露滋润这些。\OriginalTerm{米里亚}{Myria}获得了有福的份，要在白冷怀佩格，并孕育恩宠中的恩宠。犹太地已见到它的花朵，这地方正在凋零。救恩临到外邦人和异乡人，安慰丰沛地施予可怜人。妇女们应当舞蹈，说：\InnerQuote{夫人佩格，泉源之母，你天上星座的母亲。你如云彩，在炎热后给我们带来甘露，请记念你的依赖者，主母。}}\par

国王于是毫不延迟地派遣他治下的一些贤士带着礼物前往，星为他们指路。他们返回后，向当时的人们叙述了那些同样记载于金版上的事，内容如下：——\par

我们到达耶路撒冷时，这征兆连同我们的到来惊动了全体民众。他们说：\Quote{这是怎么回事？波斯的贤士们来到这里，同时还有这奇怪的天象？}犹太人的首领这样盘问我们：\Quote{这伴随你们的是什么事？你们为何而来？}我们说：\Quote{你们称为默西亚的那一位已经诞生。}他们感到困惑，不敢反驳我们。但他们对我说：\Quote{凭天主的公义，请告诉我们你们所知道的事。}我们回答他们说：\Quote{你们陷于不信；无论有没有誓言，你们都不相信我们，只随从自己轻率的意见。因为基督，至高者之子，已经诞生，祂是你们律法与会堂的颠覆者。因此，被这极美好的回答如同箭般击中，你们听到这突然临到你们的名号，心里愁苦。}然后他们商议后，催促我们接受他们的礼物，不要告诉任何人这样的事发生在他们的土地上，免得，如他们所说，\Emph{有叛乱起来反对我们}。但我们回答说：\Quote{我们为他的尊荣带来了礼物，目的是宣扬我们已知的，在他诞生时在我国发生的大事；而你们却要我们收受你们的贿赂，隐瞒那由超越诸天的神性启示给我们的事，并忽视我们本国君主的命令吗？}在向我们提出许多理由后，他们放弃了这事。当犹太王派人召见我们，与我们交谈，并就我们对他所说的话提出一些问题，我们同样行事，直到他对我们的回答大为恼怒。我们离开他，对他并不比任何普通人更加重视。\par

然后我们到了被差遣的地方，看见了母亲和孩子，星为我们指示了那位王室婴孩。我们对母亲说：\Quote{尊贵的母亲，您叫什么名字？}她说：\Quote{诸位大人，我是玛利亚。}我们对她说：\Quote{您出身何处？}她回答说：\Quote{从这白冷地区。}然后我们说：\Quote{您没有丈夫吗？}她回答：\Quote{我只是有婚约，为了婚姻的盟约，我的心思远非此事。因为我无意于此。当我对此毫不在意时，某个安息日破晓，太阳刚升起，一位天使突然向我显现，向我传报得子的喜讯。我在惊惶中呼喊：\InnerQuote{主，不要这样对我，因为我没有丈夫。}他说服我相信，我应因天主的旨意而有这个儿子。}\par

然后我们对她说：\Quote{母亲，母亲，波斯所有的神明都称您有福。您的荣耀伟大，因您被高举超过一切有名的妇女，您被证明比所有王后更女王。}\par

那孩子还坐在地上，据她说，已在第二个年头，并有几分像他的母亲。她的手很长，身体略显纤弱；肤色如熟麦；圆脸，头发束起。我们同有一位擅长写生的仆人，便将他们二人的肖像带回了本国；并由我们亲手供奉在圣殿中，上面刻有这样的铭文：\Quote{献给太阳神朱庇特，大能的天主，耶稣的君王，波斯王权所奉献。}\par

我们轮流抱起那孩子，抱在怀中，向他致敬并朝拜，献上黄金、没药和乳香，这样对他说：\Quote{我们将属于你的献给你，啊，耶稣，天上的君王。若没有你在，混乱的事物就无法被整顿。除了你的降临，天上之物无法与地上之物的结合。这样的侍奉，若只有仆人被派来，而不能有主人亲临，是无法完成的；如同国王只派总督去作战，所成就的远不及国王亲自在场。你的智慧的体系，应以这种方式对待人。}\par

那孩子在我们的抚爱和话语中跳跃欢笑。我们向母亲告别，她向我们表达了敬意，我们也以合宜的礼敬回应了她，然后回到了住宿的地方。傍晚，一个面目可怖可怕的人显现给我们，说：\Quote{快快离开，免得你们陷入罗网。}我们惊恐地问：\Quote{神圣的引领者啊，是谁胆敢谋害如此庄严的使团？}他回答说：\Quote{黑落德；但立刻起身，平安地离开。}\par

我们便匆忙离开那里；并在耶路撒冷报告了我们所见的一切。看，我们向你们讲述的关于基督的这些大事；我们看见了我们的救主基督，他被显明既是天主又是人。愿荣耀与权能归于他，至于无穷之世。阿们。\par

% structure: p0033 source=h2 target=section reason=relative-heading-rank; base=h2

\section{圣星福罗撒及其七子的受难}

本文根据吕纳的版本提供。米涅也曾引用其序言，如下：此处刊行的圣星福罗撒及其七子殉道叙事，在抄本中归属于声誉卓著的作家尤利乌斯·阿非利加努斯。或许它曾被插入其\WorkTitle{编年史}中，——欧瑟伯（\WorkTitle{教会史}，vi.31）证明该著作是极其精心撰写的，因为他在其中详述了从创世到赫利奥加巴卢斯皇帝时代的历史大事。然而，由于该著作已失传，我不敢断然肯定这篇叙事确实出自阿非利加努斯之手，尽管同时也没有理由怀疑其真实性。此外，我们根据蒙布里丘、苏里乌斯和卡杜卢斯的版本，并对照两部戈尔贝抄本和一部索邦图书馆抄本排印。这些圣人的死因与哈德良在其蒂沃利的乡间别墅所建的那座最著名的宫殿的邻近有关，据斯巴提亚努斯记载。因为当皇帝下令用一些赎罪仪式来净化这座他为享乐而建的宫殿时，司铎们趁机指控星福罗撒，声称除非星福罗撒向神祇祭献或被祭献，否则神祇不会满意；而最后这件事由哈德良执行，我们从他的许多其他事迹中知道，他极其迷信，此事大约发生在基督纪元120年，即他统治初期，实际上正如狄奥·卡西乌斯所观察，这位皇帝在那时处死了许多人。此外，所有最古老的殉道录都纪念这些殉道者，尽管日期不同。罗马殉道录连同诺特克将其庆节定于7月18日，拉巴努斯定于同月21日，乌苏阿杜斯和阿多定于6月21日。在蒂布尔提纳路，仍存在一座古老教堂的废墟，据阿林吉称（\WorkTitle{罗马地下墓穴}，iv.17），该教堂以他们的名字祝圣于天主，并仍保留着\Quote{七兄弟}的称号。我毫不怀疑它建在\WorkTitle{行传}第四节中教宗们称为\Quote{七\OriginalTerm{被暴死者}{Biothanati}}的地方，即那些被暴力致死的人，如巴罗尼乌斯在138年所评注。以上为吕纳所说；另见蒂勒蒙\WorkTitle{教会回忆录}，ii，pp.241和595；以及博朗德派\WorkTitle{圣徒行传·六月卷}，iv，p.350。\par

1. 当哈德良建造了一座宫殿，并想用那邪恶的仪式来祝圣它，开始通过祭献偶像和住在偶像中的魔鬼寻求神谕时，他们回答说：\Quote{寡妇星福罗撒和她的七个儿子，每天以祈求她的天主来伤害我们。因此，如果她和她的儿子们一起祭献，我们保证满足你的一切要求。}于是哈德良下令逮捕她和她的儿子们，并用温和的言辞劝他们同意向偶像祭献。然而，蒙福的星福罗撒回答他说：\Quote{我的丈夫格图略，连同他的兄弟阿曼丘，当他们做你手下的护民官时，为了基督的名忍受了各种刑罚，不肯同意向偶像祭献。如同优秀的运动员，他们在死亡中战胜了你的魔鬼。因为他们宁愿被斩首也不肯屈服，并为了基督的名忍受死亡：这死亡，固然在世上人中赢得了暂时的耻辱，但在天使中却赢得了永恒的尊荣和荣耀；如今他们行走在他们中间，展示他们受苦的战利品，在天上与永恒的王共享永生。}\par

2. 哈德良皇帝对圣辛佛萝莎说：\Quote{要么你和你的儿子们一起向全能的神明献祭，要么我就把你和你的儿子们一起献祭。}真福辛佛萝莎回答说：\Quote{我从哪里得到这么大的好处，居然能配得上和我的儿子们一起作为祭品献给天主？}哈德良皇帝说：\Quote{我要把你献给众神当作祭品。}真福辛佛萝莎回答说：\Quote{你的神不能把我当作祭品；但如果我为了基督、我的天主的名而被焚烧，我反而会消灭你那些魔鬼。}哈德良皇帝说：\Quote{你在这两者中选择一个：要么向我的神献祭，要么死于恶死。}真福辛佛萝莎回答说：\Quote{你认为我可以用某种恐怖来改变心意；而我是多么渴望与我的丈夫格图略一起安息，你为了基督的名处死了他。}于是哈德良皇帝下令把她带到赫丘利神庙，先在那里打她的耳光，然后把她悬吊在头发上。但当他用任何理由和任何恐怖都无法使她放弃她的善志时，他下令把一块大石头系在她的脖子上，把她扔进河里。她的兄弟欧金尼乌斯，台伯地区的首领，收起了她的尸体，并将她埋葬在同一城市的郊区。\par

3. 然后，在另一天，哈德良皇帝下令把她所有的七个儿子一起带到他面前；当他挑战他们向偶像献祭，并看到他们无论如何都不屈服于他的威胁和恐怖时，他下令在赫丘利神庙周围竖起七根木桩，并命令把他们绑在那里的刑架上。他下令：第一个克雷森斯，刺穿喉咙；第二个尤利安，刺入胸膛；第三个内梅西乌斯，刺穿心脏；第四个普里米提乌斯，伤在肚脐；第五个尤斯廷，用剑从背后刺穿；第六个斯特拉克特乌斯，伤在肋旁；第七个欧金尼乌斯，从头向下劈成两半。\par

4. 第二天，哈德良皇帝又来到赫丘利神庙，下令把他们的尸体一起运走，扔进一个深坑；祭司们给那地方起名为\Emph{七暴死者}。此后，迫害停止了一年半，在这期间，所有殉道者的圣尸都受到尊崇，并被小心地安置在为此建造的坟墓中，他们的名字记录在生命册上。此外，基督的圣殉道者，真福辛佛萝莎和她的七个儿子：克雷森斯、尤利安、内梅西乌斯、普里米提乌斯、尤斯廷、斯特拉克特乌斯和欧金尼乌斯，他们的诞辰日是在7月18日。他们的遗体安息在提布提纳路上，离城第八里程碑处，在我们的主耶稣基督的国度内，愿尊荣和荣耀归于他，直到永远。阿们。\par

% structure: p0039 source=h2 target=section reason=relative-heading-rank; base=h2

\section{断片一}

\Emph{论埃及人和迦勒底人的神话年代学}\par

埃及人确实以其自夸的古老观念，由他们的占星家以周期和万年的形式提出了一种说法；一些被认为精研此类课题的人简略地称之为太阴年；并且不亚于其他人倾向于神话，\Emph{他们认为}这与柏拉图笔下的埃及祭司向梭伦虚报的八千或九千年相符。\par

（\Emph{接着是一些其他内容：}）\par

我何必提腓尼基人的三万年，或迦勒底人的四十八万年的愚妄呢？因为犹太人，作为亚巴郎的后裔，从他们那里起源，因梅瑟的精神被教导了与真理相伴的谦逊思想和合宜的人性，通过他们现存的希伯来历史，将直到救恩圣言降临的时期——在凯撒统治时期向世界宣告——的年代数传给了我们，即五千五百年。\par

% structure: p0044 source=h2 target=section reason=relative-heading-rank; base=h2

\section{断片二}

当人类在地上增多时，天上的天使与人女结合。在一些抄本中我找到了\Quote{天主的儿子们}这句话。依我之见，圣神的意思是说：舍特的子孙被称为天主子，因为从他那里产生了直到救主本身的义人和圣祖；而加音的子孙被称为人种，因为他们里面没有神性，因他们种族的邪恶和他们本性的不平等，是一群混杂的人民，并激起了天主的愤怒。但如果认为这些是指天使，那么我们必须将他们视为那些行魔术和变戏法的人，他们教导妇女星辰的运行和天上事物的知识，藉着他们的能力，他们孕育了巨人作为他们的孩子，藉着他们，邪恶在地上达到了顶峰，直到天主因洪水判决所有活物的种族当在他们的不敬中灭亡。\par

% structure: p0046 source=h2 target=section reason=relative-heading-rank; base=h2

\section{断片三}

亚当二百三十岁时生了舍特；之后又活了七百年就死了，即第二次死亡。\par

舍特二百零五岁时生了厄诺士；因此从亚当到厄诺士出生，总共四百三十五年。\par

厄诺士一百九十岁时生了刻南。\par

刻南又一百七十岁时生了玛拉肋耳；\par

玛拉肋耳一百六十五岁时生了雅勒得；\par

雅勒得一百六十二岁时生了哈诺客；\par

哈诺客一百六十五岁时生了默突舍拉；且他中悦天主，又活了二百年之后，就不见了。\par

默突舍拉一百八十七岁时生了拉默客。\par

拉默客一百八十八岁时生了诺厄。\par

% structure: p0056 source=h2 target=section reason=relative-heading-rank; base=h2

\section{断片四 论洪水}

天主决意用洪水毁灭所有活物，曾警告人不得活过一百二十年。但不要以为这事困难，因为此后有些人活得比这更长。所指定的时期是直到洪水时的那一百年，针对当时的罪人；因为他们当时二十岁。天主指示诺厄——因他的正义而蒙主喜悦——建造一只方舟；方舟完工后，诺厄本人和他的儿子们、他的妻子和儿媳们，以及每种活物的首生者都进入方舟，以延续种族。洪水来临时，诺厄已六百岁。水退后，方舟停在阿辣辣特山上，我们知道这山在帕提亚；但有些人说它在弗里吉亚的塞莱奈，我两个地方都见过。洪水泛滥一年，然后地面干了。他们成对地从方舟出来，正如可以发现的，而不是像进去时那样按照种类区分，并蒙天主祝福。每一件事都为我们指示某种有益的事。\par

% structure: p0058 source=h2 target=section reason=relative-heading-rank; base=h2

\section{残篇五}

洪水来临时，诺厄已六百岁。因此，从亚当到诺厄和洪水，共计二千二百六十二年。\par

% structure: p0060 source=h2 target=section reason=relative-heading-rank; base=h2

\section{残篇六}

洪水以后，闪生了阿帕革沙得。\par

阿帕革沙得一百三十五岁时，在二千三百九十七年生舍拉。\par

舍拉一百三十岁时，在二千五百二十七年生厄贝尔。\par

厄贝尔一百三十四岁时，在二千六百六十一年生培肋格，因在他的日子大地被分裂而得名。\par

培肋格一百三十岁时生勒伍，又活了二百零九年而死。\par

% structure: p0066 source=h2 target=section reason=relative-heading-rank; base=h2

\section{残篇七}

在世界三千二百七十七年，亚巴郎进入应许之地客纳罕。\par

% structure: p0068 source=h2 target=section reason=relative-heading-rank; base=h2

\section{残篇八．论亚巴郎}

由此产生了\Emph{希伯来人}的名称。因为\Emph{希伯来人}一词解释为\Emph{那些迁移过来的人}，即与亚巴郎一同渡过幼发拉底河的人；它并非如有些人所想，出自前述的厄贝尔。因此，从洪水与诺厄到亚巴郎进入应许之地，共计一千零一十五年；从亚当算起，二十代，共三千二百七十七年。\par

% structure: p0070 source=h2 target=section reason=relative-heading-rank; base=h2

\section{残篇九．论亚巴郎与罗特}

当饥荒压迫客纳罕地时，亚巴郎下到埃及；因害怕他妻子的美貌会危及自己，便谎称她是自己的妹妹。但法郎——埃及人称他们的王为此名——在接到她后，将她娶了过去。他受到天主的惩罚；亚巴郎和他所有的一切都被释放，并富足地离开。在客纳罕，亚巴郎的牧人和罗特的牧人彼此相争；他们彼此同意分开，罗特因土地肥沃美好而选择住在索多玛，那里有五座城：索多玛、哈摩辣、阿德玛、责波殷、左哈尔，以及同样数目的王。他们的邻居四个叙利亚王向他们开战，其首领是以拦王革多尔老默尔。他们在盐海——今称死海——相遇。我在那里见过许多非常奇妙的事。因为那水不能维持任何活物，尸体沉入深处，而活物甚至不容易浮在水面。点燃的火把浮在水上，但熄灭后就沉下去。那里有沥青泉；它出产明矾和盐，与普通种类稍异，因为它们辛辣而透明。凡是周围所结果实，都充满浓浊的烟气。这水对使用者有治疗作用，其排干方式与其他任何水相反。若不是约但河像贝壳一样注入其中，并在很大程度上抵抗其趋势，它可能比看起来更快干涸。那里还有大量的香脂树；但一般认为，由于附近居民的不敬，它已被天主摧毁。\par

% structure: p0072 source=h2 target=section reason=relative-heading-rank; base=h2

\section{残篇十．论先祖雅各伯}

雅各伯的牧羊帐篷，一直保存在厄德萨直到罗马皇帝安敦尼的时代，后被雷电击毁。\par

雅各伯因西默盎和肋未在舍根针对当地居民所为——因他们妹妹被玷污——而不悦，便在舍根将随身携带的神像埋在一块岩石旁，在那棵奇妙的笃耨香树下；这树直到今日仍被附近居民为纪念圣祖而尊崇。雅各伯由此迁往贝特耳。在这棵笃耨香树旁有一座祭坛，当地居民在集会时献上\OriginalTerm{祭品}{ectenoe}；祭坛虽看似燃烧，却未烧尽。旁边是亚巴郎和依撒格的坟墓。有人说，曾受亚巴郎款待的一位天使的杖就插在那里。\par

% structure: p0075 source=h2 target=section reason=relative-heading-rank; base=h2

\section{残篇十一}

因此，根据这书，从亚当到若瑟的死，共二十三代，三千五百六十三年。\par

% structure: p0077 source=h2 target=section reason=relative-heading-rank; base=h2

\section{残篇十二}

因此，根据这记录，我们断言奥古格斯——阿提卡第一次洪水由此得名，他在众人死亡时获救——与梅瑟一同活在百姓出埃及的时代。\EditorialNote{有一处中断}奥古格斯之后，因洪水造成巨大破坏，阿提卡之地直到刻克洛普斯时代仍无王，共一百八十九年。然而菲洛科洛斯断言，奥古格斯、阿克泰乌斯或任何其他虚构的名称从未存在过。\EditorialNote{另一处中断}从奥古格斯到居鲁士，如同从梅瑟到他的时代，共一千二百三十五年。\par

% structure: p0079 source=h2 target=section reason=relative-heading-rank; base=h2

\section{残篇十三}

1. 直到奥利匹亚纪时代，希腊人中间并无确定的历史，此前一切事物皆混淆不清，彼此毫无一致。然而，这些\OriginalTerm{奥利匹亚纪}{Olympiads}曾由多人详细考证，因为希腊人并非按长时段，而是以四年为期来编撰其历史记录。为此，我将选取第一届奥利匹亚纪之前神话叙述中最显著的部分，并迅速加以概述。但此后的事件，至少那些重要的事件，我将按其日期与希腊事件一同处理，仔细考察希伯来人的事迹，而更粗略地涉及希腊人的事迹；我的计划如下：取希伯来历史中某一事件与希腊历史中另一事件同时发生者，以其为主线，在叙述中增减调整，记录在那一希伯来历史事件发生时，有何著名的希腊人或波斯人，或任何其他民族的显要人物活跃于世；如此，我或许可达致我所设定的目标。\par

2. 当时临到希伯来人最著名的被掳——即他们被巴比伦王\Person{拿步高}掳去之时——持续了七十年，正如\Person{耶肋米亚}所预言的。巴比伦人\Person{贝罗索斯}也提及\Person{拿步高}。被掳七十年后，\Person{居鲁士}成为波斯王，正值第五十五届奥利匹亚纪，这可以从\Person{狄奥多罗斯}的\OriginalTerm{书库}{Bibliothecae}以及\Person{塔罗斯}、\Person{卡斯托尔}的历史著作中得知，同样也见于\Person{波利比乌斯}和\Person{弗勒贡}，以及其他以奥利匹亚纪为研究对象的人。因为日期在他们所有人中间是一致的。于是，\Person{居鲁士}在他统治的第一年，即第五十五届奥利匹亚纪的第一年，藉着\Person{则鲁巴贝耳}之手实现了人民的第一次部分回归，同时还有\Person{约匝达克}之子\Person{耶叔亚}，因为七十年的时期现已届满，正如希伯来历史学家\Person{厄斯德拉}所叙述的。因此，\Person{居鲁士}统治的开始与被掳的结束相吻合。并且，根据奥利匹亚纪的推算，直至我们时代，事件之间会有类似的和谐。遵循此法，我们也将使其他叙述以同样方式彼此契合。\par

3. 但如果以\Place{阿提卡}纪年作为这些之前事件的标准，那么从\Person{奥基戈斯}（他们相信他是土著，其时代也是\Place{阿提卡}发生第一次大洪水之时，同时\Person{福洛纽斯}统治着\Place{阿尔戈斯}人，如\Person{阿库西劳斯}所述）直到第一届奥利匹亚纪（从该时期希腊人认为可以准确确定日期）总共有一千零二十年；这个数字既与上述内容相符，也将由下文所确立。因为这些事也为雅典历史学家\Person{赫拉尼库斯}和\Person{菲托科鲁斯}（他们记载阿提卡事务）以及\Person{卡斯托尔}和\Person{塔罗斯}（他们记载叙利亚事务）所记录；也为\Person{狄奥多罗斯}（他在其\OriginalTerm{书库}{Bibliothecae}中撰写通史）和\Person{亚历山大·波利希斯托}，以及我们时代一些更仔细的作者和所有阿提卡作家所记录。因此，在这1020年中出现的任何值得注意的叙述都将在适当位置给出。\par

4. 因此，根据这一著作，我们断言：\Person{奥基戈斯}（他以其名为第一次大洪水命名，并在众人灭亡时获救）生活于人民与\Person{梅瑟}一起出离埃及的时代。我们通过以下方式证明这一点。从\Person{奥基戈斯}到前述第一届奥利匹亚纪，将表明有1020年；从第一届奥利匹亚纪到第五十五届的第一年，即\Person{居鲁士}王的第一年（也就是被掳的结束），有217年。因此，从\Person{奥基戈斯}到\Person{居鲁士}是1237年。如果从被掳结束向后推算，也有1237年。这样，通过分析，从以色列在\Person{梅瑟}领导下出离埃及的第一年到第五十五届奥利匹亚纪至\Person{奥基戈斯}（他建立了\Place{厄琉西斯}）的时期相同。由此，我们为\Place{阿提卡}纪年学获得了一个更显著的起点。\par

5. 至此是关于\Person{奥基戈斯}之前时期的情况。在他之时，\Person{梅瑟}离开了埃及。我们以下列方式证明这一说法在日期上的可靠性。从\Person{梅瑟}出谷到被掳后统治的\Person{居鲁士}，共1237年。因为\Person{梅瑟}剩余的年数是40年。在他之后领导人民的\Person{若苏厄}是25年；\Person{若苏厄}之后作民长的长老们是30年；\WorkTitle{民长纪}中所记载的民长时期是490年；司祭\Person{厄里}和\Person{撒慕尔}的年数是90年；希伯来历代君王的年数是490年。\Emph{然后是被掳的七十年}，其最后一年是\Person{居鲁士}统治的第一年，正如我们已经说过的。\par

6. 从梅瑟到第一届奥林匹亚德计有1020年，从同一时期到第五十五届奥林匹亚德第一年计有1237年，在这一计数中，希腊人的计算与我们相符。在奥古吉斯之后，因洪水造成巨大破坏，阿提卡这片土地直到刻克洛普斯为止，约有189年没有国王。因为菲洛科鲁斯声称，据说接续奥古吉斯的阿克泰俄斯，或任何其他虚构的名字，从未存在过。\newline{}\Emph{再者}：因此，从奥古吉斯到居鲁士，他说，同一时期从梅瑟到同一日期，即1237年；有些希腊人也记载梅瑟生活于同一时期。例如，波勒摩在他的\WorkTitle{希腊史}第一卷中说：在福洛纽斯之子阿庇斯时代，埃及军队的一支离开埃及，定居在名叫叙利亚的巴勒斯坦，距阿拉伯不远；这些人显然就是跟随梅瑟的那些人。语法学家中最勤奋的波塞多尼奥斯之子阿皮翁，在他的\WorkTitle{驳斥犹太人}一书以及\WorkTitle{历史}第四卷中说，在阿尔戈斯国王伊那科斯时代，当阿摩西斯统治埃及时，犹太人在梅瑟领导下叛离。希罗多德也在他的第二卷书中提到这次叛离和阿摩西斯，并在某种程度上提到犹太人本身，将他们列在受割礼的人之中，称他们为巴勒斯坦的亚述人，或许是由于亚巴郎的原因。门德斯的托勒密，他从最早时期叙述埃及的历史，也对所有这一切给出了同样的记载；因此在他们中间，在年代学上总体上没有值得注意的差异。\par

7. 还应当注意，所有被希腊人因其古老而视为特别著名的传说故事，都被发现属于梅瑟之后的时期；例如他们的洪水与大火、普洛米透斯、伊娥、欧罗巴、斯巴达（播种人）、普洛塞庇娜的劫掠、他们的奥秘、他们的立法、狄俄尼索斯的事迹、珀尔修斯、阿尔戈英雄、半人马、弥诺陶洛斯、特洛伊的故事、赫拉克勒斯的劳役、赫拉克利代的回归、爱奥尼亚迁徙和奥林匹亚德。我认为特别应该记述前面提到的阿提卡王权时期，因为我打算将希腊人的历史与希伯来人的历史并列叙述。因为任何人只要从我的立场出发，就能与我同样好地做出计算。现在，在那从梅瑟和奥古吉斯延伸到第一届奥林匹亚德的1020年时期的第一年，发生了逾越节和希伯来人出离埃及，以及在阿提卡发生了奥古吉斯洪水。这是合乎理性的。因为当埃及人在天主的愤怒中被冰雹和风暴击打时，自然可以预料地球的某些部分会与他们一同受苦；特别是，可以预料雅典人应该与埃及人一同遭受这样的灾难，因为他们被认为是埃及人的一支殖民，正如忒奥蓬普斯在他的\WorkTitle{特利卡雷努斯}中以及其他人所声称的那样。中间时期被略过了，因为希腊人中没有记载什么显著事件。但94年后，据某些人所说，普洛米透斯出现了，他被神话般地传说创造了人；因为他是一个智慧的人，他把人从极度粗野的状态转变为文明。\par

% structure: p0087 source=h2 target=section reason=relative-heading-rank; base=h2

\section{断片14}

阿伽墨斯托耳之子埃斯库罗斯统治了雅典人二十三年，当时约阿堂在耶路撒冷作王。\par

我们的年表将犹大王约阿堂置于第一届奥林匹亚德之内。\par

% structure: p0090 source=h2 target=section reason=relative-heading-rank; base=h2

\section{断片15}

阿非利加诺在他的\WorkTitle{历史}第三卷中写道：\Quote{现在所记载的第一届奥林匹亚德——它实际上是第十四届——是科罗伊波斯获胜的时期；那时阿哈次在耶路撒冷作王的第一年。}在第四卷中他说：\Quote{因此，我们已经表明第一届奥林匹亚德与阿哈次作王的第一年同时发生。}\par

% structure: p0092 source=h2 target=section reason=relative-heading-rank; base=h2

\section{断片16. 论达尼尔的七十个星期}

1. 因此，这段经文如此陈述，涉及许多奇妙的事。但目前，我只谈其中涉及年代学及相关问题的事。那么，这段经文显然是论及基督的来临，祂将在七十个星期之后显现自己。因为在救主时代，或从祂开始，过犯被废除，罪恶被终止。此外，通过赦免，不义连同过犯因赎罪而被涂抹；一种永久的正义被宣讲，不同于律法所规定的；异象和预言（持续）直到若翰，并且至圣者被傅油。因为在救主来临之前，这些事尚未发生，因此只是被期待。至于数字的开端，即构成490年的七十个星期，天使教导我们从命令发出，回复并建造耶路撒冷开始计算。这事发生在波斯王阿塔薛西斯在位第二十年。因为他的酒政乃赫米雅恳求他，并得到答复说耶路撒冷应当被建造。命令发出，吩咐这些事；因为直到那时，那城是荒凉的。当居鲁士在七十年掳掠后允许所有愿意的人自由回归时，有些人在大司祭若苏厄和则鲁巴贝尔的领导下，另一些在其后的人在厄斯德拉的领导下回归了，但起初他们被阻止建造圣殿和用墙围城，借口是尚未得到命令。\par

因此，这一状况一直持续到\Person{乃赫米雅}和\Person{阿尔塔薛西斯}统治时期，以及波斯帝国统治的第一百一十五年。从耶路撒冷被攻陷算起，那是185年。那时，\Person{阿尔塔薛西斯}王下令修建城市；\Person{乃赫米雅}被派遣，监督工程，街道和城墙得以建造，正如所预言的那样。从那时算起，我们得到七十周，直到\Person{基督}时代。因为如果我们从任何其他起点而不是从这里开始计算，时期将不对应，并且会有许多奇怪的结果。因为如果我们从\Person{居鲁士}和第一次复兴开始计算七十周，就会多出一百多年；如果我们从天使向\Person{达尼尔}预言的那天开始计算，会更多；而如果我们从被掳的开始计算，则会更多得多。因为我们发现波斯帝国统治持续230年，马其顿帝国统治延续370年，从那时到\Person{提庇留·凯撒}的第十六年大约是60年。\par

因此，从\Person{阿尔塔薛西斯}开始，直到\Person{基督}的时代，按犹太人的计算，七十周得以完成。因为从\Person{乃赫米雅}（他被\Person{阿尔塔薛西斯}派遣，在波斯帝国第一百一十五年、第八十三奥林匹克第四年、\Person{阿尔塔薛西斯}本人统治第二十年，去建造耶路撒冷）直到这个日期，即第二百零二奥林匹克第二年、\Person{提庇留·凯撒}统治第十六年，共计475年，按希伯来计算法为490年，因为他们以月亮运行来度量年份；因此，很容易证明，他们的一年有354天，而太阳年有365又1/4天。因为后者超过按照月亮运行的十二个月的时间11又1/4天。因此，希腊人和犹太人每8年插入三个闰月。因为8乘以11又1/4天构成三个月。因此，475年组成59个八年周期，外加三个月。但由于每8年有三个闰月，我们得到15年\Emph{减去}数日；将这些加到475年上，总共构成七十周。\par

% structure: p0096 source=h2 target=section reason=relative-heading-rank; base=h2

\section{片段十七。论\Person{希尔卡诺斯}和\Person{安提戈诺斯}的命运，以及\Person{黑落德}、\Person{奥古斯都}、\Person{安东尼}和\Person{克娄巴特拉}，摘要}

1. \Person{屋大维·塞巴斯图}，或称\Person{奥古斯都}，这是罗马人对他的称呼，\Person{盖乌斯}的养子，从\Place{伊庇鲁斯}的\Place{阿波罗尼亚}（他受教育之地）返回罗马后，掌握了政府的最高地位。\Person{安东尼}随后获得了亚细亚及其以外地区的统治权。在他时期，犹太人控告\Person{黑落德}；但他处死了代表，并恢复了\Person{黑落德}的统治。然而后来，他与\Person{希尔卡诺斯}和其弟\Person{法萨埃卢斯}一起被驱逐，逃往\Person{安东尼}处。由于犹太人不接纳他，发生了一场激战；不久之后，他战胜了，并驱逐了返回的\Person{安提戈诺斯}。\Person{安提戈诺斯}逃往帕提亚王\Person{黑落德}那里，并借助其子\Person{帕科鲁斯}的帮助得以复位，条件是支付一千塔冷通黄金。于是\Person{黑落德}被迫逃亡，\Person{法萨埃卢斯}战死，\Person{希尔卡诺斯}被活捉交给\Person{安提戈诺斯}。\Person{安提戈诺斯}割去他的耳朵，使他不配作司铎，然后把他交给帕提亚人带去充军；因为\Person{安提戈诺斯}犹豫着是否杀他，因他是自己的亲戚。\Person{黑落德}被驱逐后，首先投奔阿拉伯王\Person{马里库斯}；但\Person{马里库斯}因惧怕帕提亚人而拒绝收留他，他就前往\Place{亚历山大港}投奔\Person{克娄巴特拉}。那是第一百八十五奥林匹克。\Person{克娄巴特拉}杀死了她的兄弟（她的统治伙伴），随后被\Person{安东尼}召到\Place{基里基雅}为自己辩护，她把统治权托付给\Person{黑落德}；\Person{黑落德}要求在他恢复自己的统治之前不被委托任何事务，于是她带他一同去见\Person{安东尼}。由于\Person{安东尼}爱上了这位女王，他们派遣\Person{黑落德}去\Place{罗马}见\Person{屋大维·奥古斯都}；\Person{奥古斯都}因\Person{黑落德}的父亲\Person{安提帕特}以及\Person{黑落德}本人的缘故，又因\Person{安提戈诺斯}凭借帕提亚人的帮助自立为王，于是命令\Place{巴勒斯坦}和\Place{叙利亚}的将领恢复\Person{黑落德}的统治。\Person{黑落德}与\Person{索修斯}联手，长期与\Person{安提戈诺斯}作战，进行了多次交锋。那时，\Person{黑落德}的兄弟\Person{若瑟夫}也在指挥中阵亡。\Person{黑落德}来到\Person{安东尼}……\par

2. 他们围攻\Person{安提戈诺斯}三年，然后将他活捉交给\Person{安东尼}。\Person{安东尼}本人也宣布\Person{黑落德}为王，并赐给他\Place{希普斯}、\Place{加达拉}、\Place{迦萨}、\Place{约培}、\Place{安特东}等城市，以及阿拉伯的一部分、\Place{特拉可尼}、\Place{奥兰尼提斯}、\Place{萨西亚}、\Place{哥兰尼提斯}；此外还有叙利亚的总督职位。\Person{黑落德}被元老院和\Person{屋大维·奥古斯都}宣布为犹太人的王，统治了34年。\Person{安东尼}在即将远征帕提亚时，杀了犹太王\Person{安提戈诺斯}，并将阿拉伯赐给\Person{克娄巴特拉}；他进入帕提亚境内后遭受惨败，损失了大部分军队。这是在第一百八十六奥林匹克。\Person{屋大维·奥古斯都}率领意大利和整个西部的军队对抗\Person{安东尼}，\Person{安东尼}因惧怕（由于在帕提亚的失败）以及出于对\Person{克娄巴特拉}的爱，拒绝返回罗马。\Person{安东尼}率领亚细亚的军队迎战。然而\Person{黑落德}像一个精明的人，一个侍奉权贵的人，他发出两套书信，并将军队派往海上，命令将领们观察事态的发展。当胜负已定，\Person{安东尼}在两次海战失败后与\Person{克娄巴特拉}一同逃往埃及时，携带书信的人将那些他们一直为\Person{安东尼}秘密保存的书信交给了\Person{奥古斯都}。于是\Person{黑落德}……\par

3. 克娄帕特拉将自己封闭在一座陵墓中，以毒蛇为死亡工具自尽。当时奥古斯都在赫利奥斯与塞勒涅逃往底比斯地区途中俘获了克娄帕特拉的儿子们。尼科波利斯建在阿克提翁对面，并设立了称为阿克提亚竞技会的赛会。亚历山大城被攻占后，科尔内利乌斯·加卢斯被派往埃及担任首任总督，他摧毁了那些拒绝服从的埃及城市。至此，拉戈斯王朝统治结束；波斯势力覆亡后，马其顿帝国共持续了二百九十八年。从马其顿帝国建立到托勒密王朝覆亡——末代女王克娄帕特拉统治时期——整个时期至此完成，该事件发生在罗马君主国与帝国的第十一年，以及第一百八十七届奥林匹克周期的第四年。从亚当算起，总共五千四百七十二年。\par

4. 亚历山大城被攻占后，第一百八十八届奥林匹克周期开始。黑落德重建了加比尼人的古城撒玛黎雅，并称之为塞巴斯特；他又将港口斯特拉托塔扩建为城市，以同名命名为凯撒勒雅，并在两城各建了一座神庙以尊崇屋大维。之后，他在吕底亚平原建立了安提帕特里斯，以其父命名，并将他从塞巴斯特地区驱逐的人民安置在那里。他还建立了其他城市；他对犹太人严厉，而对其他民族则极为和蔼。\par

第一百八十九届奥林匹克周期，该周期在闰日那一年（即三月朔日前第六天，亦即二月二十四日）与安提约基雅纪元的第二十四年对应，从而确定了该年的准确界限。\par

% structure: p0102 source=h2 target=section reason=relative-heading-rank; base=h2

\section{片段十八　论与救主受难及祂赋予生命的复活相关的情况}

1. 至于祂的各样事迹、祂在灵魂与身体上施行的医治、祂教导的奥秘以及死者复活，这些已由祂的门徒和在我们之前的宗徒们最具权威地陈明了。整个大地笼罩着一片极其可怕的黑暗；岩石因地震而崩裂，犹太地区及其他许多地方坍塌。塔洛斯在其\WorkTitle{历史}第三卷中称这黑暗为日蚀，但在我看来并无道理。因为希伯来人在月中的第十四天庆祝逾越节，而我们的救主受难恰在逾越节前一日；但日蚀只发生在月亮运行到太阳之下时。并且它只能发生在新月第一天与旧月最后一天之间的间隔，即月朔之时：那么，当月亮几乎与太阳正相对时，怎么会发生日蚀呢？不过，让这种观点过去吧；让它获得多数人的认同；让这个世界的征兆被视为一次日蚀——如同其他仅对眼睛而言的征兆。弗莱贡记载，在提庇留·凯撒时期，月圆之时，从第六时辰到第九时辰有一次完整的日蚀——显然就是我们所说的那一次。但日蚀与地震、岩石崩裂、死者复活以及整个宇宙如此巨大的动荡有何共同之处？确实，如此事件在漫长时期内并无记录。然而，这是由于天主造成的黑暗，因为主那时正受苦难。计算显示，达尼尔所记的七十个星期在此时期完成了。\par

2. 此外，从阿塔薛西斯算起，按照犹太人的计算，到基督时期共七十个星期。因为从乃赫米雅（由阿塔薛西斯派往耶路撒冷为人民安居）——约在波斯帝国第一百二十年、阿塔薛西斯本人第二十年、第八十三届奥林匹克周期第四年——直到现在，即第一百零二届奥林匹克周期第二年、提庇留·凯撒统治第十六年，共四百七十五年，这相当于四百九十个希伯来年，因为他们以每月二十九日半的阴历月计算，这很容易解释：太阳年周期为三百六十五日又四分之一日，而十二个月阴历月少十一又四分之一日。因此，希腊人和犹太人每八年插入三个闰月。因为八乘以十一又四分之一日等于三个月。所以，四百七十五年包含五十九个八年周期再加三个月：这样，每八年加上三个闰月，得到十五年，这十五年与四百七十五年相加即为七十个星期。现在，请不要认为我们不谙天文计算，因为我们未经进一步讨论就确定一年为三百六十五又四分之一日。这并非出于对真相的无知，而是出于精确的研究才简要陈述了我们的观点。但对于那些力求详尽探究一切的人，也以纲要的方式呈现如下。\par

3. 每年通常由三百六十五天组成；一日一夜被分成十九份，我们还有五份。说一年有三百六十五又四分之一天，并且有五个十九分之一……对四百七十五年来说，有六又四分之一天。此外，根据精确计算，我们发现阴历月有二十九日半……而这一点只差很少时间。从阿塔薛西斯统治第二十年（如希伯来人的\BibleBook{厄斯德拉}{}所记，按希腊算法为第八十届奥林匹克周期第四年）到提庇留·凯撒第十六年（第一百零二届奥林匹克周期第二年），总共是上述的四百七十五年，按希伯来体系计算为四百九十年，正如先前所述，即七十个星期，加俾额尔向达尼尔宣布的基督来临时间正是以此衡量的。如果有人以为所加的十五年希伯来年会导致十年的误差，那至少没有引入任何无法解释的问题。而且，我们为凑足总数所假定加上的一个半星期，也解决了关于这十五年的疑问，消除了时间上的困难；预言通常以某种象征形式呈现，这一点是相当明显的。\par

4. 那么，尽我们所能，我认为我们正确地理解了圣经，尤其是考虑到前文关于异象的部分似乎简短地陈述了整个事件，其开头的话是：\Quote{在贝耳沙匝王第三年，}\BibleRef{达 8:1} 他在那里预言了波斯帝国被希腊人推翻，这两个帝国在预言中分别以公绵羊和公山羊的形象为象征。\BibleRef{达 8:13} \Quote{祭献，}他说，\Quote{要被废除，圣所要被荒废，以致被践踏；这些事要在二千三百日内决定。}\BibleRef{达 8:13} 因为如果我们以一天为一个月，正如在预言的其他地方天数被当作年数，并且在不同地方有不同的用法，把这段时间同样按上述方法换算为希伯来月，我们就会发现从耶路撒冷被攻陷到阿塔薛西斯王第二十年，这段时间完全吻合。因为共有一百八十五年，还要加上一年——乃赫米雅修建城墙的那一年。因此，在一百八十六年中，我们找到了二百三十个希伯来月，因为八年中还有三个闰月。再从阿塔薛西斯——在他统治时期，建造耶路撒冷的命令发出——算起，有七十个星期。不过，这些事我们在\WorkTitle{论星期与本预言}一书中已单独且更精确地讨论过。但我感到惊奇，犹太人竟否认主已经来临，马尔基翁的追随者也拒绝承认祂的来临曾在预言中被预言，而圣经如此公开地将这事显明在我们眼前。\Emph{另外还有：}那么，从亚当和创造到主降临的时期是五千五百三十一年，从这个时代到第二百五十届奥林匹亚竞技会有一百九十二年，如上所述。\par

% structure: p0107 source=h2 target=section reason=relative-heading-rank; base=h2

\section{残篇十九}

因为我们既知道那些话语的分量，又不忽视信德的恩宠，便感谢圣父，祂赐给我们这些受造物耶稣基督，万有的救主，我们的主；愿光荣与威严归于祂，偕同圣神，直到永远。\par

\SourceRecord{Translator unknown.；Ante-Nicene Fathers；Vol. 6.；Edited by Alexander Roberts, James Donaldson, and A. Cleveland Coxe.；Buffalo, NY: Christian Literature Publishing Co.,；1886.；<http://www.newadvent.org/fathers/0614.htm>.}
